\scanpage{10}\mspage{II}
% L1–4
Soit $V$ un \deleted{\unsure{espace}}plan vectoriel sur $\mathbb F_3$. La droite projective $P(V)$ a quatre éléments i.e. $V$ \deleted{\unsure{a}}\inserted{contient} quatre droites vectorielles -- donnons nous une paire \inserted{$\unsure{\tau|T}$} de droites de $V$ $(D_i)_{i\in\tau}$. On a donc
\[
V=\mathop{\deleted{\prod}\mkern-15mu\bigoplus}_{i\in\tau}D_i\simeq\prod_{i\in\tau}D_i
\tag*{(1)}
\]
% L6–8
Considérons la paire complémentaire $\tau'=P(V)\smallsetminus\tau$, \deleted{\unsure{elle}}\unsure{on} a donc aussi
\[
V=\bigoplus_{i\in\tau'}D_i\simeq\prod_{i\in\tau'}D_i
\tag*{(2)}
\]
% L9–15
Notons les relations
\[
\left\{\begin{aligned}
\tau'&\simeq\bigwedge_{i\in\tau}D_i^*\\
\tau&\simeq\bigwedge_{i\in\tau'}D_i^*
\end{aligned}\right.
\tag*{(3)}
\]
l'application
\[
\prod_{i\in\tau}D_i^*\longrightarrow\tau'
\]
qui définit $\bigwedge_{i\in\tau}^{i}D_i^*\unsure{\mathord{-}|\to}\tau'$ étant
\[
(u_i)_{i\in\tau}\longrightarrow\mathbb F_3\mathbin{\cdot}
\left(\sum_{i\in\tau}u_i\right)\quad.
\tag*{(4)}
\]
% L16–21
Si $P$ est un $V$-torseur, i.e. un \deleted{\unsure{espace|plan}}\inserted{plan} affine sous $\mathbb F_3$ admettant $V$ comme groupe de translations, \unsure{alors|donc} les formules (1) et (2) définissent respectivement deux décompositions de $P$ en produits
\[
P\unsure{=|\simeq}\prod_{i\in\tau}P_i\unsure{\mathord{-}|\mathord{.}},\qquad
P\simeq\prod_{i\in\tau'}P_i\unsure{.|'}
\tag*{(5)}
\]
% archive sheet number: 9

